% =============================================================================
%  PAES Studio: presentación de repaso (plantilla Beamer 16:9)
%  Lista para Overleaf: Nuevo proyecto → Subir proyecto → compilar con pdfLaTeX.
%
%  Estética minimalista con la paleta Zinc de la app: sin degradados, sin
%  símbolos de navegación, líneas finas y alternativas con su tecla A–E.
%  Comandos:
%    \alternativa{A}{texto}       alternativa
%    \alternativaok{C}{texto}     alternativa correcta (para la diapositiva de solución)
%    \paso{n}{texto}      paso numerado del desarrollo
%    \ejebarra{Eje}{0.42}{42 %}   barra proporcional de frecuencia
%  Entornos: cajaclave, cajatrampa, lineatiempo (filas «fecha & hecho \\»).
%
%  El generador scripts/pipeline/generate_latex_summary.py reemplaza los bloques
%  delimitados por %%<NOMBRE>%% … %%</NOMBRE>%%.
% =============================================================================
\documentclass[aspectratio=169,11pt]{beamer}

\usepackage{lmodern}        % antes de fontenc/babel: no depende de las fuentes EC
\usepackage[T1]{fontenc}
\usepackage[utf8]{inputenc}
\usepackage[spanish,es-noshorthands,es-tabla]{babel}
\spanishplainpercent        % \% sin espaciado automático (evita choques en modo matemático)
\usepackage{textcomp}
\usepackage{amsmath,amssymb}
\usepackage{booktabs}
\usepackage{array}
\usepackage[most]{tcolorbox}

% ── Paleta Zinc ──────────────────────────────────────────────────────────────
\definecolor{zinc950}{HTML}{09090B}
\definecolor{zinc900}{HTML}{18181B}
\definecolor{zinc700}{HTML}{3F3F46}
\definecolor{zinc500}{HTML}{71717A}
\definecolor{zinc300}{HTML}{D4D4D8}
\definecolor{zinc200}{HTML}{E4E4E7}
\definecolor{zinc50}{HTML}{FAFAFA}
\definecolor{esmeralda}{HTML}{047857}
\definecolor{rosa}{HTML}{BE123C}

\newcommand{\sen}{\operatorname{sen}}

% ── Tema mínimo ──────────────────────────────────────────────────────────────
\usefonttheme{professionalfonts}
\setbeamertemplate{navigation symbols}{}
\setbeamercolor{background canvas}{bg=white}
\setbeamercolor{normal text}{fg=zinc900}
\setbeamercolor{structure}{fg=zinc900}
\setbeamercolor{alerted text}{fg=zinc950}
\setbeamercolor{item}{fg=zinc500}
\setbeamerfont{frametitle}{series=\bfseries,size=\large}
\setbeamertemplate{itemize items}{\raisebox{0.15ex}{\color{zinc500}\scriptsize$\bullet$}}
\setbeamertemplate{enumerate items}{\color{zinc500}\insertenumlabel.}
\setbeamersize{text margin left=1.6cm,text margin right=1.6cm}

\makeatletter
\setbeamertemplate{frametitle}{%
  \vspace*{0.9em}%
  {\usebeamerfont{frametitle}\color{zinc950}\insertframetitle}\par
  \ifx\insertframesubtitle\@empty\else{\small\color{zinc500}\insertframesubtitle}\par\fi
  \vspace{0.25em}{\color{zinc200}\rule{\textwidth}{0.4pt}}%
}
\makeatother
\setbeamertemplate{footline}{%
  \hspace*{1.6cm}{\color{zinc500}\tiny\insertshorttitle}\hfill
  {\color{zinc500}\tiny\ttfamily\insertframenumber/\inserttotalframenumber}\hspace*{1.6cm}\par\vspace*{0.7em}%
}
\setbeamertemplate{title page}{%
  \vfill
  {\color{zinc500}\small PAES Studio}\par\vspace{0.6em}
  {\usebeamerfont{title}\bfseries\Huge\color{zinc950}\inserttitle}\par\vspace{0.5em}
  {\large\color{zinc700}\insertsubtitle}\par\vspace{1.2em}
  {\color{zinc300}\rule{0.35\textwidth}{0.4pt}}\par\vspace{0.6em}
  {\small\color{zinc500}\insertdate}\par
  \vfill
}

% ── Componentes ──────────────────────────────────────────────────────────────
\newtcbox{\tecla}{on line, boxrule=0.4pt, arc=1.5pt, colback=zinc50, colframe=zinc300,
  left=3pt, right=3pt, top=1pt, bottom=1pt, fontupper=\ttfamily\footnotesize\color{zinc700}}
\newtcbox{\teclaok}{on line, boxrule=0.4pt, arc=1.5pt, colback=esmeralda, colframe=esmeralda,
  left=3pt, right=3pt, top=1pt, bottom=1pt, fontupper=\ttfamily\footnotesize\bfseries\color{white}}
\newtcolorbox{opcion}[1][]{enhanced, boxrule=0.4pt, arc=2pt, colback=white, colframe=zinc300,
  left=6pt, right=6pt, top=3pt, bottom=3pt, before skip=3pt, after skip=3pt, #1}
\newcommand{\alternativa}[2]{\begin{opcion}\tecla{#1}\hspace{0.6em}#2\end{opcion}}
\newcommand{\alternativaok}[2]{\begin{opcion}[colframe=esmeralda!60, colback=esmeralda!5]\teclaok{#1}\hspace{0.6em}#2\end{opcion}}
\newcommand{\paso}[2]{\par\noindent\makebox[1.4em][l]{\color{zinc500}\ttfamily\small #1}\parbox[t]{\dimexpr\linewidth-1.4em}{#2}\par\vspace{0.45em}}
\newtcolorbox{cajaclave}[1][]{enhanced, boxrule=0.4pt, arc=2pt, colback=zinc50, colframe=zinc300,
  borderline west={2pt}{0pt}{zinc900}, left=7pt, right=7pt, top=3pt, bottom=4pt,
  fonttitle=\bfseries\small, coltitle=zinc900, colbacktitle=zinc50, titlerule=0pt, title={#1}}
\newtcolorbox{cajatrampa}[1][Trampas frecuentes]{enhanced, boxrule=0.4pt, arc=2pt, colback=white, colframe=rosa!35,
  borderline west={2pt}{0pt}{rosa}, left=7pt, right=7pt, top=3pt, bottom=4pt,
  fonttitle=\bfseries\small, coltitle=zinc900, colbacktitle=white, titlerule=0pt, title={#1}}
\newenvironment{lineatiempo}
  {\footnotesize\renewcommand{\arraystretch}{1.15}\begin{tabular}{@{}>{\leavevmode\bfseries\color{zinc700}}p{0.14\textwidth}p{0.74\textwidth}@{}}}
  {\end{tabular}}
\newlength{\anchobarra}
\newcommand{\ejebarra}[3]{%
  \par\noindent\makebox[0.34\linewidth][l]{\small #1}%
  \setlength{\anchobarra}{0.5\linewidth}%
  \makebox[\anchobarra][l]{{\color{zinc200}\rule[0.2ex]{\anchobarra}{1.1ex}}\hspace*{-\anchobarra}{\color{zinc900}\rule[0.2ex]{#2\anchobarra}{1.1ex}}}%
  \hspace{0.8em}{\small\ttfamily #3}\par\vspace{0.5em}}

%%<META>%%
\title[Repaso PAES: Ciencias]{Repaso PAES}
\subtitle{Ciencias}
\date{8 de noviembre de 2026}
%%</META>%%

\begin{document}

\begin{frame}[plain]
  \titlepage
\end{frame}

%%<CONTENIDO>%%
\begin{frame}{Contenidos}

  \framesubtitle{4 ejes, en el orden del temario}

  \begin{enumerate}\small

    \item \textbf{Biología}: Organización, estructura y actividad celular, Sistema nervioso, Reproducción, sexualidad e ITS, Ciclo celular, meiosis y manipulación genética, Evolución: evidencias y selección natural, Fotosíntesis, respiración y flujo de energía

    \item \textbf{Física}: Ondas electromagnéticas, luz y óptica, Movimiento: MRU y MRUA, Fuerzas y leyes de Newton, Sistema solar y universo, Tectónica de placas y sismos, Circuitos eléctricos y consumo de energía

    \item \textbf{Química}: Materia y estructura atómica, Química orgánica, Reacciones químicas y estequiometría, Disoluciones y concentración

    \item \textbf{Habilidades científicas}: Habilidades de pensamiento científico

  \end{enumerate}

\end{frame}

\begin{frame}{Qué evalúa la prueba}

  \framesubtitle{Distribución por eje en 21 ensayos (1474 preguntas)}

  \vspace{0.8em}

  \ejebarra{Biología}{0.338}{34 \%}

  \ejebarra{Física}{0.338}{34 \%}

  \ejebarra{Química}{0.324}{32 \%}

\end{frame}

\section{Biología}

\begin{frame}[plain]
  \vfill
  {\color{zinc500}\small Eje}\par\vspace{0.3em}{\Large\bfseries Biología}\par\vspace{0.6em}
  {\color{zinc700} Organización celular, procesos y funciones biológicas (sistema nervioso y reproducción), herencia y evolución, organismo y ambiente.}\par\vspace{1em}
  {\small\color{zinc500} Organización, estructura y actividad celular, Sistema nervioso, Reproducción, sexualidad e ITS, Ciclo celular, meiosis y manipulación genética, Evolución: evidencias y selección natural, Fotosíntesis, respiración y flujo de energía}
  \vfill
\end{frame}

\begin{frame}[allowframebreaks]{Organización, estructura y actividad celular}
  \framesubtitle{Teoría esencial}
  \begin{itemize}\small
    \item Célula vegetal: pared celular, cloroplastos y vacuola grande.
    \item Osmosis: el agua pasa hacia el medio más concentrado en solutos.
    \item Los virus no son células.
  \end{itemize}
  \begin{cajaclave}[Organelos y función]\footnotesize
Núcleo: material genético. Mitocondria: respiración celular (ATP). Cloroplasto: fotosíntesis (solo vegetales). Ribosoma: síntesis de proteínas. Retículo endoplasmático y aparato de Golgi: síntesis, modificación y transporte. Lisosoma: digestión celular.
  \end{cajaclave}
\begin{center}\footnotesize
\begin{tabular}{@{}>{\raggedright\arraybackslash}p{\dimexpr(\linewidth-4\tabcolsep)/3\relax}>{\raggedright\arraybackslash}p{\dimexpr(\linewidth-4\tabcolsep)/3\relax}>{\raggedright\arraybackslash}p{\dimexpr(\linewidth-4\tabcolsep)/3\relax}@{}}
\toprule
 & \textbf{Procarionte} & \textbf{Eucarionte} \\
\midrule
Núcleo & no (nucleoide) & sí \\
Organelos con membrana & no & sí \\
Ejemplos & bacterias & animales, plantas, hongos \\
\bottomrule
\end{tabular}
\end{center}
\end{frame}

\begin{frame}{Desafío}
  \framesubtitle{Organización, estructura y actividad celular}
  {\small Un glóbulo rojo se coloca en agua destilada. ¿Qué ocurrirá con la célula?}
  \vspace{0.4em}
  \alternativa{A}{Se deshidrata, porque el agua sale de ella.}
  \alternativa{B}{Se hincha y puede estallar, porque el agua entra en ella.}
  \alternativa{C}{No cambia, porque la membrana es impermeable al agua.}
  \alternativa{D}{Pierde sus sales hacia el agua por transporte activo.}
\end{frame}

\begin{frame}[allowframebreaks]{Solución}
  \framesubtitle{El agua destilada es hipotónica respecto del interior celular.}
  {\small
  \paso{1}{El interior del glóbulo rojo tiene más solutos que el agua destilada.}
  \paso{2}{Por osmosis, el agua entra hacia el medio más concentrado: la célula se hincha (hemólisis).}
  }
  \alternativaok{B}{Se hincha y puede estallar, porque el agua entra en ella.}
  \begin{cajatrampa}\footnotesize
    \begin{itemize}
    \item Atribuir cloroplastos o pared celular a células animales.
    \item Pensar que en la osmosis se mueve el soluto (se mueve el agua).
    \end{itemize}
  \end{cajatrampa}
\end{frame}

\begin{frame}[allowframebreaks]{Sistema nervioso}
  \framesubtitle{Teoría esencial}
  \begin{itemize}\small
    \item En la sinapsis química el mensaje pasa con neurotransmisores.
    \item El reflejo es rápido porque no lo procesa el cerebro.
    \item Las drogas de abuso actúan sobre los neurotransmisores.
  \end{itemize}
  \begin{cajaclave}[Arco reflejo]\footnotesize
Receptor $\rightarrow$ neurona sensitiva (aferente) $\rightarrow$ médula espinal (interneurona) $\rightarrow$ neurona motora (eferente) $\rightarrow$ efector (músculo o glándula).
  \end{cajaclave}
\end{frame}

\begin{frame}{Desafío}
  \framesubtitle{Sistema nervioso}
  {\small Una persona toca una superficie caliente y retira la mano antes de sentir dolor. ¿Qué explica esa rapidez?}
  \vspace{0.4em}
  \alternativa{A}{El estímulo llega primero a la corteza cerebral.}
  \alternativa{B}{La respuesta se integra en la médula espinal, en un arco reflejo.}
  \alternativa{C}{Las hormonas actúan más rápido que los nervios.}
  \alternativa{D}{La piel mueve los músculos por sí sola.}
\end{frame}

\begin{frame}[allowframebreaks]{Solución}
  \framesubtitle{Es un reflejo: la médula integra la respuesta sin esperar al cerebro.}
  {\small
  \paso{1}{La neurona sensitiva lleva el estímulo a la médula.}
  \paso{2}{Desde ahí, una neurona motora activa el músculo; el dolor se percibe después.}
  }
  \alternativaok{B}{La respuesta se integra en la médula espinal, en un arco reflejo.}
  \begin{cajatrampa}\footnotesize
    \begin{itemize}
    \item Pensar que todo reflejo pasa por el cerebro.
    \item Confundir dendrita (recibe) con axón (conduce el impulso).
    \end{itemize}
  \end{cajatrampa}
\end{frame}

\begin{frame}[allowframebreaks]{Reproducción, sexualidad e ITS}
  \framesubtitle{Teoría esencial}
  \begin{itemize}\small
    \item Ovulación: liberación del ovocito, cerca del día 14.
    \item Solo los métodos de barrera protegen de las ITS.
    \item El VIH es un virus; la gonorrea y la clamidia, bacterias.
  \end{itemize}
  \begin{cajaclave}[Ciclo de 28 días]\footnotesize
Días 1–5: menstruación. Días 6–13: fase folicular, el endometrio se engrosa. Día 14 aprox.: ovulación. Días 15–28: fase lútea; si no hay fecundación, bajan las hormonas y comienza un nuevo ciclo.
  \end{cajaclave}
\end{frame}

\begin{frame}{Desafío}
  \framesubtitle{Reproducción, sexualidad e ITS}
  {\small ¿Cuál de los siguientes métodos previene a la vez el embarazo y las infecciones de transmisión sexual?}
  \vspace{0.4em}
  \alternativa{A}{Píldora anticonceptiva.}
  \alternativa{B}{Método del calendario.}
  \alternativa{C}{Preservativo.}
  \alternativa{D}{Ligadura de trompas.}
\end{frame}

\begin{frame}[allowframebreaks]{Solución}
  \framesubtitle{Solo una barrera física impide el contacto con fluidos y patógenos.}
  {\small
  \paso{1}{La píldora y la ligadura evitan el embarazo, pero no las ITS.}
  \paso{2}{El calendario no protege de ninguna de las dos de forma confiable.}
  }
  \alternativaok{C}{Preservativo.}
  \begin{cajatrampa}\footnotesize
    \begin{itemize}
    \item Creer que los anticonceptivos hormonales protegen de las ITS.
    \item Ubicar la ovulación al inicio del ciclo.
    \end{itemize}
  \end{cajatrampa}
\end{frame}

\begin{frame}[allowframebreaks]{Ciclo celular, meiosis y manipulación genética}
  \framesubtitle{Teoría esencial}
  \begin{itemize}\small
    \item La replicación del ADN ocurre en la fase $S$, no en la mitosis.
    \item La variabilidad viene de la meiosis, no de la mitosis.
    \item Un transgénico porta un gen de otra especie.
  \end{itemize}
  \begin{cajaclave}[Mitosis y meiosis]\footnotesize
Mitosis: $2n\rightarrow$ 2 células $2n$ idénticas. Meiosis: $2n\rightarrow$ 4 células $n$ distintas.

Humanos: $2n=46$ en células somáticas y $n=23$ en gametos. La cromatina se compacta en cromosomas visibles durante la división.
  \end{cajaclave}
\begin{center}\footnotesize
\begin{tabular}{@{}>{\raggedright\arraybackslash}p{\dimexpr(\linewidth-4\tabcolsep)/3\relax}>{\raggedright\arraybackslash}p{\dimexpr(\linewidth-4\tabcolsep)/3\relax}>{\raggedright\arraybackslash}p{\dimexpr(\linewidth-4\tabcolsep)/3\relax}@{}}
\toprule
 & \textbf{Mitosis} & \textbf{Meiosis} \\
\midrule
Células hijas & 2 & 4 \\
Cromosomas & igual a la madre ($2n$) & la mitad ($n$) \\
Función & crecimiento y reparación & formar gametos \\
\bottomrule
\end{tabular}
\end{center}
\end{frame}

\begin{frame}{Desafío}
  \framesubtitle{Ciclo celular, meiosis y manipulación genética}
  {\small Una célula humana somática con 46 cromosomas completa una mitosis. ¿Cuántos cromosomas tiene cada célula hija?}
  \vspace{0.4em}
  \alternativa{A}{23}
  \alternativa{B}{46}
  \alternativa{C}{69}
  \alternativa{D}{92}
\end{frame}

\begin{frame}[allowframebreaks]{Solución}
  \framesubtitle{La mitosis conserva el número de cromosomas.}
  {\small
  \paso{1}{Antes de dividirse, cada cromosoma se duplicó en la fase $S$.}
  \paso{2}{En la anafase se separan las cromátidas y cada hija recibe 46 cromosomas.}
  }
  \alternativaok{B}{46}
  \begin{cajatrampa}\footnotesize
    \begin{itemize}
    \item Atribuir a la mitosis la variabilidad genética.
    \item Confundir la duplicación del ADN (fase $S$) con la separación de cromátidas (anafase).
    \end{itemize}
  \end{cajatrampa}
\end{frame}

\begin{frame}[allowframebreaks]{Evolución: evidencias y selección natural}
  \framesubtitle{Teoría esencial}
  \begin{itemize}\small
    \item Evolucionan las poblaciones, no los individuos.
    \item Las mutaciones ocurren al azar; la selección no es azarosa.
    \item Lamarck (caracteres adquiridos) fue descartado.
  \end{itemize}
  \begin{cajaclave}[Selección natural]\footnotesize
Variabilidad $\rightarrow$ presión del ambiente $\rightarrow$ supervivencia y reproducción diferencial $\rightarrow$ cambio en la frecuencia de los rasgos en la población.
  \end{cajaclave}
\end{frame}

\begin{frame}{Desafío}
  \framesubtitle{Evolución: evidencias y selección natural}
  {\small Tras usar un antibiótico por años, en un hospital predominan bacterias resistentes. ¿Cuál es la mejor explicación?}
  \vspace{0.4em}
  \alternativa{A}{El antibiótico produjo mutaciones dirigidas a la resistencia.}
  \alternativa{B}{Las bacterias resistentes ya existían y sobrevivieron y se reprodujeron más.}
  \alternativa{C}{Las bacterias aprendieron a resistir el antibiótico.}
  \alternativa{D}{El antibiótico perdió su efecto por el paso del tiempo.}
\end{frame}

\begin{frame}[allowframebreaks]{Solución}
  \framesubtitle{Es un ejemplo clásico de selección natural.}
  {\small
  \paso{1}{La resistencia surgió por variación previa (mutaciones al azar).}
  \paso{2}{El antibiótico eliminó a las sensibles y las resistentes dejaron más descendencia.}
  }
  \alternativaok{B}{Las bacterias resistentes ya existían y sobrevivieron y se reprodujeron más.}
  \begin{cajatrampa}\footnotesize
    \begin{itemize}
    \item Pensar que el individuo cambia «para adaptarse».
    \item Explicar con herencia de caracteres adquiridos.
    \end{itemize}
  \end{cajatrampa}
\end{frame}

\begin{frame}[allowframebreaks]{Fotosíntesis, respiración y flujo de energía}
  \framesubtitle{Teoría esencial}
  \begin{itemize}\small
    \item Las plantas también respiran, día y noche.
    \item El $\mathrm{O_2}$ liberado proviene del agua.
    \item La energía fluye; la materia se recicla.
  \end{itemize}
  \begin{cajaclave}[Ecuaciones globales]\footnotesize
\[6\,\mathrm{CO_2}+6\,\mathrm{H_2O}\xrightarrow{\text{luz}}\mathrm{C_6H_{12}O_6}+6\,\mathrm{O_2}\]

Respiración: la reacción inversa, que libera la energía de la glucosa como ATP.
  \end{cajaclave}
\end{frame}

\begin{frame}{Desafío}
  \framesubtitle{Fotosíntesis, respiración y flujo de energía}
  {\small En un experimento, una planta acuática se ilumina con luz cada vez más intensa. ¿Qué es lo más probable que ocurra con la liberación de burbujas de oxígeno?}
  \vspace{0.4em}
  \alternativa{A}{Disminuye siempre.}
  \alternativa{B}{Aumenta hasta un punto y luego se estabiliza.}
  \alternativa{C}{No cambia.}
  \alternativa{D}{Aumenta sin límite.}
\end{frame}

\begin{frame}[allowframebreaks]{Solución}
  \framesubtitle{La luz es un factor limitante solo hasta que otro factor limita el proceso.}
  {\small
  \paso{1}{Más luz acelera la etapa luminosa y la liberación de $\mathrm{O_2}$.}
  \paso{2}{Luego el $\mathrm{CO_2}$ o la temperatura limitan y la tasa se estabiliza.}
  }
  \alternativaok{B}{Aumenta hasta un punto y luego se estabiliza.}
  \begin{cajatrampa}\footnotesize
    \begin{itemize}
    \item Creer que las plantas no respiran.
    \item Decir que la energía se recicla en el ecosistema.
    \end{itemize}
  \end{cajatrampa}
\end{frame}

\section{Física}

\begin{frame}[plain]
  \vfill
  {\color{zinc500}\small Eje}\par\vspace{0.3em}{\Large\bfseries Física}\par\vspace{0.6em}
  {\color{zinc700} Ondas electromagnéticas, mecánica (movimiento, fuerzas y universo), energía y Tierra, y electricidad.}\par\vspace{1em}
  {\small\color{zinc500} Ondas electromagnéticas, luz y óptica, Movimiento: MRU y MRUA, Fuerzas y leyes de Newton, Sistema solar y universo, Tectónica de placas y sismos, Circuitos eléctricos y consumo de energía}
  \vfill
\end{frame}

\begin{frame}[allowframebreaks]{Ondas electromagnéticas, luz y óptica}
  \framesubtitle{Teoría esencial}
  \begin{itemize}\small
    \item Todas las ondas electromagnéticas viajan a la misma rapidez en el vacío.
    \item Los ángulos de reflexión y refracción se miden respecto de la normal.
    \item El arcoíris se explica por refracción y dispersión.
  \end{itemize}
  \begin{cajaclave}[Relaciones]\footnotesize
\begin{gather*}
c=\lambda\,f \\
T=\frac{1}{f} \\
\theta_i=\theta_r \\
n=\frac{c}{v}
\end{gather*}

Más frecuencia $\Leftrightarrow$ menor longitud de onda $\Leftrightarrow$ más energía (radio $<$ microondas $<$ infrarrojo $<$ visible $<$ UV $<$ X $<$ gamma).
  \end{cajaclave}
\end{frame}

\begin{frame}{Desafío}
  \framesubtitle{Ondas electromagnéticas, luz y óptica}
  {\small Una radio transmite con una frecuencia de $100\ \text{MHz}=1\cdot 10^8$ Hz. ¿Cuál es la longitud de onda de esa señal en el aire?}
  \vspace{0.4em}
  \alternativa{A}{0,3 m}
  \alternativa{B}{1 m}
  \alternativa{C}{3 m}
  \alternativa{D}{30 m}
\end{frame}

\begin{frame}[allowframebreaks]{Solución}
  \framesubtitle{Las ondas de radio son electromagnéticas: viajan a unos $3\cdot 10^8$ m/s.}
  {\small
  \paso{1}{$\lambda=\dfrac{c}{f}=\dfrac{3\cdot 10^8}{1\cdot 10^8}$.}
  \paso{2}{$\lambda=3$ m.}
  }
  \alternativaok{C}{3 m}
  \begin{cajatrampa}\footnotesize
    \begin{itemize}
    \item Pensar que la frecuencia cambia al pasar de un medio a otro (cambian la rapidez y la longitud de onda).
    \item Creer que los rayos X o la luz necesitan un medio para propagarse.
    \end{itemize}
  \end{cajatrampa}
\end{frame}

\begin{frame}[allowframebreaks]{Movimiento: MRU y MRUA}
  \framesubtitle{Teoría esencial}
  \begin{itemize}\small
    \item Pendiente de $x$–$t$ = velocidad.
    \item Área bajo $v$–$t$ = desplazamiento.
    \item 72 km/h = 20 m/s (divide por 3,6).
  \end{itemize}
  \begin{cajaclave}[Ecuaciones]\footnotesize
\begin{gather*}
\text{MRU: } x=x_0+vt \\
\text{MRUA: } v=v_0+at,\quad x=x_0+v_0t+\tfrac{1}{2}at^2
\end{gather*}
  \end{cajaclave}
\end{frame}

\begin{frame}{Desafío}
  \framesubtitle{Movimiento: MRU y MRUA}
  {\small Un ciclista recorre 18 km en 45 minutos. ¿Cuál es su rapidez media?}
  \vspace{0.4em}
  \alternativa{A}{0,4 km/h}
  \alternativa{B}{2,5 km/h}
  \alternativa{C}{24 km/h}
  \alternativa{D}{40 km/h}
\end{frame}

\begin{frame}[allowframebreaks]{Solución}
  \framesubtitle{Rapidez media = distancia / tiempo, con unidades coherentes.}
  {\small
  \paso{1}{45 min = 0,75 h.}
  \paso{2}{$v=\dfrac{18}{0{,}75}=24$ km/h.}
  }
  \alternativaok{C}{24 km/h}
  \begin{cajatrampa}\footnotesize
    \begin{itemize}
    \item Mezclar km/h con m/s.
    \item Olvidar el factor $\frac{1}{2}$ en $\frac{1}{2}at^2$.
    \end{itemize}
  \end{cajatrampa}
\end{frame}

\begin{frame}[allowframebreaks]{Fuerzas y leyes de Newton}
  \framesubtitle{Teoría esencial}
  \begin{itemize}\small
    \item Acción y reacción actúan sobre cuerpos distintos: no se anulan.
    \item Velocidad constante $\Rightarrow$ fuerza neta nula.
    \item En la Luna cambia el peso, no la masa.
  \end{itemize}
  \begin{cajaclave}[Fórmulas]\footnotesize
\begin{gather*}
F_{\text{neta}}=m\,a \\
P=m\,g \\
F_{\text{elástica}}=k\,x \\
f_r=\mu\,N
\end{gather*}
  \end{cajaclave}
\end{frame}

\begin{frame}{Desafío}
  \framesubtitle{Fuerzas y leyes de Newton}
  {\small Un auto avanza en línea recta a 60 km/h constantes. ¿Qué se puede afirmar de la fuerza neta sobre él?}
  \vspace{0.4em}
  \alternativa{A}{Es nula.}
  \alternativa{B}{Apunta en el sentido del movimiento.}
  \alternativa{C}{Apunta en sentido contrario al movimiento.}
  \alternativa{D}{Es igual a su peso.}
\end{frame}

\begin{frame}[allowframebreaks]{Solución}
  \framesubtitle{Velocidad constante significa aceleración cero.}
  {\small
  \paso{1}{Si $a=0$, por la segunda ley $F_{\text{neta}}=m\cdot 0=0$.}
  \paso{2}{La fuerza del motor se equilibra con los roces.}
  }
  \alternativaok{A}{Es nula.}
  \begin{cajatrampa}\footnotesize
    \begin{itemize}
    \item Creer que se necesita una fuerza para mantener la velocidad constante.
    \item Sumar acción y reacción como si actuaran sobre el mismo cuerpo.
    \end{itemize}
  \end{cajatrampa}
\end{frame}

\begin{frame}[allowframebreaks]{Sistema solar y universo}
  \framesubtitle{Teoría esencial}
  \begin{itemize}\small
    \item Las lunas de Júpiter mostraron que no todo gira en torno a la Tierra.
    \item Un planeta es más rápido en su perihelio que en su afelio.
    \item Corrimiento al rojo: la galaxia se aleja.
  \end{itemize}
  \begin{cajaclave}[Leyes de Kepler]\footnotesize
1.ª: los planetas giran en órbitas elípticas con el Sol en un foco. 2.ª: la línea Sol–planeta barre áreas iguales en tiempos iguales (el planeta va más rápido cerca del Sol). 3.ª: el cuadrado del período es proporcional al cubo del semieje mayor, $T^2\propto a^3$.
  \end{cajaclave}
\end{frame}

\begin{frame}{Desafío}
  \framesubtitle{Sistema solar y universo}
  {\small ¿Qué observación de Galileo con su telescopio contradijo el modelo geocéntrico?}
  \vspace{0.4em}
  \alternativa{A}{Que la Tierra es redonda.}
  \alternativa{B}{Que Júpiter tiene lunas que giran a su alrededor.}
  \alternativa{C}{Que las órbitas son elípticas.}
  \alternativa{D}{Que el universo se expande.}
\end{frame}

\begin{frame}[allowframebreaks]{Solución}
  \framesubtitle{El modelo geocéntrico afirmaba que todos los astros giraban en torno a la Tierra.}
  {\small
  \paso{1}{Las lunas de Júpiter giran en torno a otro cuerpo, no a la Tierra.}
  \paso{2}{Las órbitas elípticas son de Kepler, y la expansión del universo es del siglo XX.}
  }
  \alternativaok{B}{Que Júpiter tiene lunas que giran a su alrededor.}
  \begin{cajatrampa}\footnotesize
    \begin{itemize}
    \item Atribuir a Copérnico las órbitas elípticas, que son de Kepler.
    \item Pensar que el Big Bang fue una explosión en un punto del espacio y no una expansión del espacio mismo.
    \end{itemize}
  \end{cajatrampa}
\end{frame}

\begin{frame}[allowframebreaks]{Tectónica de placas y sismos}
  \framesubtitle{Teoría esencial}
  \begin{itemize}\small
    \item Chile está sobre un borde convergente: la placa de Nazca subduce bajo la Sudamericana.
    \item Magnitud: una por sismo; intensidad: distinta en cada lugar.
    \item Las ondas sísmicas se propagan por el interior de la Tierra.
  \end{itemize}
  \begin{cajaclave}[Bordes de placa]\footnotesize
Convergente: las placas chocan; subducción, sismos, volcanes y cordilleras (Chile).

Divergente: se separan; dorsales y nuevo suelo oceánico. Transformante: se deslizan de lado; sismos (falla de San Andrés).
  \end{cajaclave}
\end{frame}

\begin{frame}{Desafío}
  \framesubtitle{Tectónica de placas y sismos}
  {\small ¿Por qué Chile tiene tantos sismos y volcanes?}
  \vspace{0.4em}
  \alternativa{A}{Porque está en el centro de una placa tectónica.}
  \alternativa{B}{Porque la placa de Nazca subduce bajo la placa Sudamericana.}
  \alternativa{C}{Porque las placas se separan bajo el país.}
  \alternativa{D}{Porque el núcleo terrestre está más cerca de la superficie en Chile.}
\end{frame}

\begin{frame}[allowframebreaks]{Solución}
  \framesubtitle{Chile está en un borde convergente.}
  {\small
  \paso{1}{En la subducción, la placa oceánica se hunde bajo la continental.}
  \paso{2}{La fricción y la fusión del material generan sismos y volcanismo.}
  }
  \alternativaok{B}{Porque la placa de Nazca subduce bajo la placa Sudamericana.}
  \begin{cajatrampa}\footnotesize
    \begin{itemize}
    \item Confundir epicentro con hipocentro.
    \item Pensar que la intensidad de un sismo es igual en todas las ciudades.
    \end{itemize}
  \end{cajatrampa}
\end{frame}

\begin{frame}[allowframebreaks]{Circuitos eléctricos y consumo de energía}
  \framesubtitle{Teoría esencial}
  \begin{itemize}\small
    \item Las casas conectan los artefactos en paralelo (mismo voltaje).
    \item 1 kWh es la energía de un artefacto de 1000 W funcionando una hora.
    \item La resistencia equivalente en paralelo es menor que la menor de ellas.
  \end{itemize}
  \begin{cajaclave}[Fórmulas]\footnotesize
\begin{gather*}
V=I\,R \\
P=V\,I \\
E=P\,t \\
R_{\text{serie}}=R_1+R_2 \\
\frac{1}{R_{\text{paralelo}}}=\frac{1}{R_1}+\frac{1}{R_2}
\end{gather*}
  \end{cajaclave}
\begin{center}\footnotesize
\begin{tabular}{@{}>{\raggedright\arraybackslash}p{\dimexpr(\linewidth-4\tabcolsep)/3\relax}>{\raggedright\arraybackslash}p{\dimexpr(\linewidth-4\tabcolsep)/3\relax}>{\raggedright\arraybackslash}p{\dimexpr(\linewidth-4\tabcolsep)/3\relax}@{}}
\toprule
 & \textbf{Serie} & \textbf{Paralelo} \\
\midrule
Corriente & igual en todas & se reparte \\
Voltaje & se reparte & igual en todas \\
$R$ equivalente & $R_1+R_2$ & menor que cada una \\
\bottomrule
\end{tabular}
\end{center}
\end{frame}

\begin{frame}{Desafío}
  \framesubtitle{Circuitos eléctricos y consumo de energía}
  {\small Un hervidor funciona con 220 V y por él circulan 10 A. ¿Cuál es su potencia?}
  \vspace{0.4em}
  \alternativa{A}{22 W}
  \alternativa{B}{230 W}
  \alternativa{C}{2200 W}
  \alternativa{D}{22 000 W}
\end{frame}

\begin{frame}[allowframebreaks]{Solución}
  \framesubtitle{$P=VI$.}
  {\small
  \paso{1}{$P=220\cdot 10=2200$ W.}
  }
  \alternativaok{C}{2200 W}
  \begin{cajatrampa}\footnotesize
    \begin{itemize}
    \item Sumar resistencias en paralelo como si estuvieran en serie.
    \item Confundir potencia (W) con energía consumida (kWh).
    \end{itemize}
  \end{cajatrampa}
\end{frame}

\section{Química}

\begin{frame}[plain]
  \vfill
  {\color{zinc500}\small Eje}\par\vspace{0.3em}{\Large\bfseries Química}\par\vspace{0.6em}
  {\color{zinc700} Estructura atómica, química orgánica, reacciones y estequiometría, disoluciones.}\par\vspace{1em}
  {\small\color{zinc500} Materia y estructura atómica, Química orgánica, Reacciones químicas y estequiometría, Disoluciones y concentración}
  \vfill
\end{frame}

\begin{frame}[allowframebreaks]{Materia y estructura atómica}
  \framesubtitle{Teoría esencial}
  \begin{itemize}\small
    \item Un catión perdió electrones; un anión los ganó.
    \item La destilación separa líquidos por su punto de ebullición; la filtración, un sólido de un líquido.
    \item El agua potable es una mezcla; el agua destilada, una sustancia pura.
  \end{itemize}
  \begin{cajaclave}[Relaciones]\footnotesize
\begin{gather*}
n=A-Z \\
e^-=Z-\text{carga} \\
\rho=\frac{m}{V}
\end{gather*}

Modelos: Dalton $\rightarrow$ Thomson $\rightarrow$ Rutherford $\rightarrow$ Bohr.
  \end{cajaclave}
\end{frame}

\begin{frame}{Desafío}
  \framesubtitle{Materia y estructura atómica}
  {\small El ion ${}^{32}_{16}\mathrm{S}^{2-}$ tiene:}
  \vspace{0.4em}
  \alternativa{A}{16 protones y 16 electrones.}
  \alternativa{B}{16 protones y 18 electrones.}
  \alternativa{C}{18 protones y 16 electrones.}
  \alternativa{D}{32 protones y 18 electrones.}
\end{frame}

\begin{frame}[allowframebreaks]{Solución}
  \framesubtitle{La carga $2-$ indica que ganó dos electrones.}
  {\small
  \paso{1}{Protones: $Z=16$.}
  \paso{2}{Electrones: $16+2=18$.}
  }
  \alternativaok{B}{16 protones y 18 electrones.}
  \begin{cajatrampa}\footnotesize
    \begin{itemize}
    \item Tomar $A$ como número de neutrones.
    \item Usar filtración para separar una mezcla homogénea.
    \end{itemize}
  \end{cajatrampa}
\end{frame}

\begin{frame}[allowframebreaks]{Química orgánica}
  \framesubtitle{Teoría esencial}
  \begin{itemize}\small
    \item Cada doble enlace aporta un enlace pi ($\pi$).
    \item Carbono con doble enlace: $sp^2$.
    \item Alcanos: solo enlaces simples ($\mathrm{C_nH_{2n+2}}$).
  \end{itemize}
  \begin{cajaclave}[Grupos funcionales]\footnotesize
Alcohol: $\mathrm{R{-}OH}$. Aldehído: $\mathrm{R{-}CHO}$. Cetona: $\mathrm{R{-}CO{-}R'}$. Ácido: $\mathrm{R{-}COOH}$. Éter: $\mathrm{R{-}O{-}R'}$. Éster: $\mathrm{R{-}COO{-}R'}$. Amina: $\mathrm{R{-}NH_2}$.
  \end{cajaclave}
\end{frame}

\begin{frame}{Desafío}
  \framesubtitle{Química orgánica}
  {\small ¿Cuál es la fórmula molecular del propano, un alcano de 3 carbonos?}
  \vspace{0.4em}
  \alternativa{A}{$\mathrm{C_3H_4}$}
  \alternativa{B}{$\mathrm{C_3H_6}$}
  \alternativa{C}{$\mathrm{C_3H_8}$}
  \alternativa{D}{$\mathrm{C_3H_{10}}$}
\end{frame}

\begin{frame}[allowframebreaks]{Solución}
  \framesubtitle{Alcanos: $\mathrm{C_nH_{2n+2}}$.}
  {\small
  \paso{1}{$n=3$: $2\cdot 3+2=8$ hidrógenos.}
  \paso{2}{$\mathrm{C_3H_8}$.}
  }
  \alternativaok{C}{$\mathrm{C_3H_8}$}
  \begin{cajatrampa}\footnotesize
    \begin{itemize}
    \item Confundir aldehído (carbonilo en el extremo) con cetona (carbonilo interno).
    \item Contar mal los enlaces pi.
    \end{itemize}
  \end{cajatrampa}
\end{frame}

\begin{frame}[allowframebreaks]{Reacciones químicas y estequiometría}
  \framesubtitle{Teoría esencial}
  \begin{itemize}\small
    \item Balancea antes de calcular.
    \item Los coeficientes relacionan moles, no gramos.
    \item Masa total de reactantes = masa total de productos.
  \end{itemize}
  \begin{cajaclave}[Relaciones]\footnotesize
\begin{gather*}
n=\frac{m}{M} \\
N=n\cdot 6{,}02\cdot 10^{23} \\
\%\,\text{elemento}=\frac{\text{masa del elemento}}{\text{masa molar}}\cdot 100
\end{gather*}

Leyes ponderales: conservación de la masa (Lavoisier), proporciones definidas (Proust) y múltiples (Dalton). La fórmula empírica es la proporción mínima entre átomos; la molecular, un múltiplo de ella.
  \end{cajaclave}
\end{frame}

\begin{frame}{Desafío}
  \framesubtitle{Reacciones químicas y estequiometría}
  {\small En la reacción balanceada $\mathrm{N_2}+3\,\mathrm{H_2}\rightarrow 2\,\mathrm{NH_3}$, ¿cuántos moles de $\mathrm{NH_3}$ se forman a partir de 6 mol de $\mathrm{H_2}$ con exceso de $\mathrm{N_2}$?}
  \vspace{0.4em}
  \alternativa{A}{2 mol}
  \alternativa{B}{3 mol}
  \alternativa{C}{4 mol}
  \alternativa{D}{6 mol}
\end{frame}

\begin{frame}[allowframebreaks]{Solución}
  \framesubtitle{La proporción entre $\mathrm{H_2}$ y $\mathrm{NH_3}$ es $3:2$.}
  {\small
  \paso{1}{$6\ \text{mol H}_2\cdot\dfrac{2\ \text{mol NH}_3}{3\ \text{mol H}_2}=4$ mol de $\mathrm{NH_3}$.}
  }
  \alternativaok{C}{4 mol}
  \begin{cajatrampa}\footnotesize
    \begin{itemize}
    \item Usar las proporciones de coeficientes directamente con gramos.
    \item Olvidar balancear la ecuación.
    \end{itemize}
  \end{cajatrampa}
\end{frame}

\begin{frame}[allowframebreaks]{Disoluciones y concentración}
  \framesubtitle{Teoría esencial}
  \begin{itemize}\small
    \item Diluir no cambia los moles de soluto.
    \item Molaridad: por litro de \textbf{disolución}; molalidad: por kilogramo de \textbf{disolvente}.
    \item La solubilidad de los gases baja al subir la temperatura.
  \end{itemize}
  \begin{cajaclave}[Fórmulas]\footnotesize
\begin{gather*}
M=\frac{n_{\text{soluto}}}{V_{\text{disolución}}\,(\text{L})} \\
m=\frac{n_{\text{soluto}}}{\text{kg disolvente}} \\
\%\,m/v=\frac{g\ \text{soluto}}{mL\ \text{disolución}}\cdot 100 \\
C_1V_1=C_2V_2
\end{gather*}
  \end{cajaclave}
\end{frame}

\begin{frame}{Desafío}
  \framesubtitle{Disoluciones y concentración}
  {\small Se diluyen 100 mL de una disolución 2 mol/L hasta 400 mL. ¿Cuál es la nueva concentración?}
  \vspace{0.4em}
  \alternativa{A}{0,5 mol/L}
  \alternativa{B}{1 mol/L}
  \alternativa{C}{4 mol/L}
  \alternativa{D}{8 mol/L}
\end{frame}

\begin{frame}[allowframebreaks]{Solución}
  \framesubtitle{Los moles de soluto se conservan: $C_1V_1=C_2V_2$.}
  {\small
  \paso{1}{$2\cdot 100=C_2\cdot 400$.}
  \paso{2}{$C_2=0{,}5$ mol/L.}
  }
  \alternativaok{A}{0,5 mol/L}
  \begin{cajatrampa}\footnotesize
    \begin{itemize}
    \item Usar mililitros en vez de litros en la molaridad.
    \item Confundir masa de disolvente con masa de disolución.
    \end{itemize}
  \end{cajatrampa}
\end{frame}

\section{Habilidades científicas}

\begin{frame}[plain]
  \vfill
  {\color{zinc500}\small Eje}\par\vspace{0.3em}{\Large\bfseries Habilidades científicas}\par\vspace{0.6em}
  {\color{zinc700} Pensamiento científico: variables, hipótesis, procedimientos y conclusiones.}\par\vspace{1em}
  {\small\color{zinc500} Habilidades de pensamiento científico}
  \vfill
\end{frame}

\begin{frame}[allowframebreaks]{Habilidades de pensamiento científico}
  \framesubtitle{Teoría esencial}
  \begin{itemize}\small
    \item Independiente: lo que yo cambio. Dependiente: lo que mido.
    \item Una conclusión solo puede afirmar lo que los datos muestran.
    \item El grupo control permite comparar el efecto de la variable.
  \end{itemize}
  \begin{cajaclave}[Esquema de una investigación]\footnotesize
Observación $\rightarrow$ pregunta de investigación $\rightarrow$ hipótesis $\rightarrow$ diseño experimental (variables y grupo control) $\rightarrow$ resultados $\rightarrow$ conclusión.
  \end{cajaclave}
\end{frame}

\begin{frame}{Desafío}
  \framesubtitle{Habilidades de pensamiento científico}
  {\small Una estudiante riega cuatro plantas iguales con distinta cantidad de agua y mide su altura después de un mes. ¿Cuál es la variable dependiente?}
  \vspace{0.4em}
  \alternativa{A}{La cantidad de agua.}
  \alternativa{B}{El tipo de planta.}
  \alternativa{C}{La altura de las plantas.}
  \alternativa{D}{El tiempo de un mes.}
\end{frame}

\begin{frame}[allowframebreaks]{Solución}
  \framesubtitle{La dependiente es la que se mide como respuesta.}
  {\small
  \paso{1}{Manipula: cantidad de agua (independiente).}
  \paso{2}{Mide: altura (dependiente). Tipo de planta y tiempo son controladas.}
  }
  \alternativaok{C}{La altura de las plantas.}
  \begin{cajatrampa}\footnotesize
    \begin{itemize}
    \item Confundir la variable que se mide con la que se manipula.
    \item Elegir una conclusión que generaliza más allá de los datos.
    \end{itemize}
  \end{cajatrampa}
\end{frame}
%%</CONTENIDO>%%

\end{document}
